\section{Capítulo 11 -- Testes de API}

Até aqui, cada capítulo testou o TaskFlow à mão: rodar o servidor, enviar requisições com o \texttt{curl}, o REST Client ou o Postman, e conferir as respostas com os próprios olhos. Esse cuidado encontrou defeitos, mas não se repete sozinho: a cada mudança no código, tudo precisaria ser conferido de novo. Neste capítulo você vai escrever testes automáticos com o Vitest: testes de unidade, para as funções que não dependem de nada; testes de integração, que usam um banco de dados de verdade; testes das rotas, sem precisar do servidor; e um teste de contrato, que confere as respostas contra o documento OpenAPI do Capítulo 10. No final, um único comando confere, em poucos segundos, o que antes levava uma tarde.

\subsection{11.1 Por que testar automaticamente}

Um teste automático é um programa que usa o seu código e confere o resultado. Ele não substitui o raciocínio de quem programa, mas traz três coisas que o teste manual não traz:

\begin{itemize}
  \item \textbf{Repetição sem custo.} Os 34 testes deste capítulo rodam em cerca de 10 segundos. Rodá-los a cada mudança é barato; repetir à mão o roteiro de testes dos Capítulos 6 a 10 não é.
  \item \textbf{Proteção contra regressões.} Uma \textbf{regressão} é um defeito em algo que já funcionava, causado por uma mudança em outro lugar. É o tipo mais comum de defeito, e o mais difícil de perceber à mão, porque ninguém volta a testar o que ``já estava pronto''.
  \item \textbf{Coragem para mudar.} Com os testes passando, uma reorganização do código, como a camada de serviço do Capítulo 7 ou o tratador de erros do Capítulo 8, pode ser feita sem medo: se algo quebrar, um teste avisa.
\end{itemize}

Os testes também servem de documentação. Um teste chamado ``não deixa mexer na tarefa de outra pessoa'' explica uma regra da API melhor que um comentário, e, ao contrário do comentário, avisa quando deixa de ser verdade.

\subsection{11.2 Unidade, integração e ponta a ponta}

Os testes costumam ser divididos pelo tamanho do pedaço do sistema que exercitam:

\begin{center}
\begin{tabular}{@{}P{0.2\linewidth}P{0.36\linewidth}P{0.34\linewidth}@{}}
\toprule
\rowcolor{headergray}
\tblheader{Tipo} & \tblheader{O que exercita} & \tblheader{No TaskFlow} \\
\midrule
unidade & uma função ou módulo isolado, sem banco nem rede & os esquemas do Zod, o \texttt{problema}, os \textit{tokens} \\
\rowcolor{rowgray} integração & várias partes juntas, com dependências reais, como o banco & os serviços e as rotas, com um banco de testes \\
ponta a ponta & o sistema inteiro, do jeito que o cliente o usa & a coleção do Postman rodada pelo Newman (Capítulo 10) \\
\bottomrule
\end{tabular}
\end{center}

Os testes de unidade são os mais rápidos e apontam o defeito com precisão, mas não veem os problemas que surgem quando as partes se juntam, como uma consulta do Prisma que não faz o que se esperava. Os de ponta a ponta veem tudo, mas são lentos, precisam do servidor no ar e, quando falham, não dizem onde está o defeito. Por isso a recomendação clássica, a \textbf{pirâmide de testes}, é ter muitos testes de unidade, uma boa quantidade de testes de integração e poucos de ponta a ponta. Em APIs, em que boa parte da lógica está nas consultas ao banco, os testes de integração costumam ser os que mais valem.

\subsection{11.3 Instalando o Vitest}

O \textbf{Vitest} é um \textit{framework} de testes para JavaScript e TypeScript. Ele roda o TypeScript sem configuração extra, entende os módulos modernos (\texttt{import}) e é rápido. A forma de escrever os testes é praticamente a mesma do \textbf{Jest}, o \textit{framework} mais antigo e ainda muito usado: \texttt{describe}, \texttt{it} e \texttt{expect} funcionam igual nos dois, e quem aprende um sabe usar o outro.

\begin{codigo}
\begin{Verbatim}[breaklines=true, breaksymbolleft={}, fontsize=\footnotesize]
$ npm install -D vitest@5 @vitest/coverage-v8@5 @types/node@24 ajv
\end{Verbatim}
\end{codigo}

O \texttt{@vitest/coverage-v8} mede a cobertura dos testes (Seção 11.11), e o \texttt{ajv} será usado no teste de contrato (Seção 11.10). O \texttt{@types/node@24} está na lista por um motivo prático: o \texttt{create-next-app} instala os tipos do Node.js 20, e o Vitest 5 exige os do 22 ou mais novo. Sem ele, o \texttt{npm} recusa a instalação com um erro \texttt{ERESOLVE}, ``conflito de dependências''. Instalar os tipos da versão do Node.js que você usa resolve. A configuração fica em um arquivo na raiz do projeto:

\begin{codigo}
\begin{Verbatim}[breaklines=true, breaksymbolleft={}, fontsize=\footnotesize]
// vitest.config.mts -- a configuração dos testes
import { loadEnv } from "vite";
import { defineConfig } from "vitest/config";

export default defineConfig({
  resolve: {
    // o mesmo atalho "@/" do tsconfig.json
    alias: { "@": import.meta.dirname },
  },
  test: {
    // carrega o .env e, por cima dele, o .env.test
    env: loadEnv("test", import.meta.dirname, ""),
    // prepara o banco dos testes uma vez, antes de tudo
    globalSetup: "./testes/preparar-banco.ts",
    // os arquivos dividem o mesmo banco: um de cada vez
    fileParallelism: false,
    coverage: {
      include: ["lib/**", "app/api/**"],
      exclude: ["lib/openapi.ts"],
    },
  },
});
\end{Verbatim}
\end{codigo}

O arquivo usa a extensão \texttt{.mts}, de ``módulo TypeScript'': o \texttt{package.json} do Next.js não declara o projeto como módulo, e com \texttt{.ts} o Vitest mostra um aviso a cada execução. Cada opção tem um motivo:

\begin{itemize}
  \item o \texttt{alias} ensina ao Vitest o atalho \texttt{@/}, que o Next.js lê do \texttt{tsconfig.json}; sem ele, nenhum \texttt{import} do projeto funcionaria nos testes;
  \item o \texttt{loadEnv("test", ...)} carrega as variáveis do \texttt{.env} e, por cima delas, as do \texttt{.env.test}, assunto da próxima seção;
  \item o \texttt{globalSetup} roda um arquivo uma única vez, antes de todos os testes;
  \item o \texttt{fileParallelism: false} faz os arquivos de teste rodarem um de cada vez. Por padrão, o Vitest roda vários ao mesmo tempo, o que é ótimo para testes de unidade, mas aqui todos usam o mesmo banco, e um apagaria os dados do outro no meio do caminho;
  \item o \texttt{coverage} diz quais arquivos entram na medição de cobertura. O documento OpenAPI fica de fora: ele é uma descrição, não uma lógica a testar.
\end{itemize}

Por fim, dois \textit{scripts} no \texttt{package.json}:

\begin{codigo}
\begin{Verbatim}[breaklines=true, breaksymbolleft={}, fontsize=\footnotesize]
"scripts": {
  "dev": "next dev",
  "build": "next build",
  "start": "next start",
  "test": "vitest",
  "cobertura": "vitest run --coverage"
}
\end{Verbatim}
\end{codigo}

O \texttt{npm test} abre o Vitest no \textbf{modo de observação} (\textit{watch}): ele roda os testes e fica esperando; a cada arquivo salvo, roda de novo só os testes afetados. É o jeito mais confortável de trabalhar. Em servidores de integração contínua, onde não há ninguém olhando, o Vitest percebe a variável \texttt{CI} e roda uma única vez. O \texttt{npm run cobertura} roda uma vez e mede a cobertura.

\subsection{11.4 Um banco só para os testes}

Os testes de integração criam e apagam dados à vontade: cada teste começa com o banco vazio, para não depender do que outro teste deixou. Por isso eles precisam de um banco \textbf{só deles}, e nunca o de desenvolvimento, onde estão os dados do \textit{seed} e os seus testes manuais, e muito menos o de produção. Crie um segundo banco, como fez no Capítulo 3: no Neon, um \textit{branch} chamado \texttt{testes}, ou, no computador, outra instância do \texttt{prisma dev}, com \texttt{npx prisma dev --name taskflow-testes}. Depois, coloque a \textit{connection string} dele em um arquivo próprio:

\begin{codigo}
\begin{Verbatim}[breaklines=true, breaksymbolleft={}, fontsize=\footnotesize]
# .env.test -- o banco dos testes (os testes APAGAM os dados)
DATABASE_URL="postgresql://...o banco dos testes..."
DIRECT_URL="postgresql://...o banco dos testes..."
\end{Verbatim}
\end{codigo}

O \texttt{.gitignore} do Next.js já protege todos os arquivos que começam com \texttt{.env}. O \texttt{JWT\_SECRET} não precisa ser repetido: o \texttt{loadEnv} carrega o \texttt{.env} primeiro, e o \texttt{.env.test} só troca o que ele define. O banco dos testes começa vazio, sem as tabelas. Quem as cria é o arquivo do \texttt{globalSetup}, que também confere se o banco é mesmo outro:

\begin{codigo}
\begin{Verbatim}[breaklines=true, breaksymbolleft={}, fontsize=\footnotesize]
// testes/preparar-banco.ts -- roda uma vez, antes de todos os
// testes: confere o banco dos testes e aplica as migrações
import { execSync } from "node:child_process";
import { existsSync, readFileSync } from "node:fs";
import { parse } from "dotenv";

export default function prepararBanco() {
  if (!existsSync(".env.test")) {
    throw new Error("Crie o .env.test com o banco dos testes.");
  }
  const desenvolvimento = parse(readFileSync(".env"));
  const testes = parse(readFileSync(".env.test"));

  // Os testes apagam dados: nunca podem rodar no banco de
  // desenvolvimento, e muito menos no de produção
  if (
    !testes.DATABASE_URL ||
    testes.DATABASE_URL === desenvolvimento.DATABASE_URL
  ) {
    throw new Error(
      "O .env.test precisa apontar para OUTRO banco.",
    );
  }

  execSync("npx prisma migrate deploy", {
    stdio: "inherit",
    env: { ...process.env, ...testes },
  });
}
\end{Verbatim}
\end{codigo}

A comparação é uma trava de segurança: se o \texttt{.env.test} faltar, ou apontar para o mesmo banco do \texttt{.env}, os testes nem começam. Um descuido aqui apagaria todos os dados de desenvolvimento na primeira execução; em um projeto real, com o \texttt{.env} apontando para produção, o estrago seria muito pior. O \texttt{execSync} roda um comando no terminal e espera ele terminar. O \texttt{prisma migrate deploy} aplica as migrações que ainda não foram aplicadas, sem criar novas e sem perguntar nada, o comando certo para bancos que não são o de desenvolvimento. As variáveis do \texttt{.env.test} vão no \texttt{env}, por cima das do sistema; o \texttt{dotenv} do \texttt{prisma7.config.ts} não sobrescreve variáveis que já existem, então o Prisma usa o banco dos testes. O comando só é executado na primeira vez em que há migrações novas; nas outras, ele responde ``No pending migrations to apply''.

Algumas funções vão ser usadas por vários arquivos de teste:

\begin{codigo}
\begin{Verbatim}[breaklines=true, breaksymbolleft={}, fontsize=\footnotesize]
// testes/apoio.ts -- funções que vários testes usam
import { db } from "@/lib/db";
import { cadastrarUsuario } from "@/lib/servicos/usuarios";

// Cada teste começa com o banco vazio
export async function limparBanco() {
  await db.task.deleteMany();
  await db.user.deleteMany();
}

let contador = 0;

// Cadastra uma pessoa com um e-mail que ainda não existe
export async function criarPessoa(nome = "Pessoa de Teste") {
  contador++;
  const usuario = await cadastrarUsuario({
    name: nome,
    email: `pessoa${contador}@teste.com`,
    password: "senha12345",
  });
  if (!usuario) throw new Error("E-mail repetido no teste.");
  return usuario;
}
\end{Verbatim}
\end{codigo}

A \texttt{limparBanco} apaga as tarefas antes dos usuários, por causa da chave estrangeira. A \texttt{criarPessoa} usa o próprio serviço de cadastro, com um contador para que cada pessoa tenha um e-mail diferente.

\subsection{11.5 O primeiro teste}

Os testes ficam na pasta \texttt{testes}, em arquivos terminados em \texttt{.test.ts}, que é o padrão que o Vitest procura. Comece pelos esquemas do Zod, que são funções puras: recebem dados e devolvem dados, sem banco nem rede.

\begin{codigo}
\begin{Verbatim}[breaklines=true, breaksymbolleft={}, fontsize=\footnotesize]
// testes/unidade/esquemas.test.ts -- o que a API aceita
import { describe, expect, it } from "vitest";
import {
  esquemaFiltros,
  esquemaNovaTarefa,
} from "@/lib/esquemas/tarefas";
import { esquemaCadastro } from "@/lib/esquemas/auth";

describe("esquemaNovaTarefa", () => {
  it("aceita só o título e tira os espaços das pontas", () => {
    const dados = esquemaNovaTarefa.parse({
      title: "  Estudar ",
    });
    expect(dados).toEqual({ title: "Estudar" });
  });

  it("recusa o título vazio, com a mensagem certa", () => {
    const resultado = esquemaNovaTarefa.safeParse({ title: "" });
    expect(resultado.success).toBe(false);
    expect(resultado.error?.issues[0].message).toBe(
      "O título é obrigatório.",
    );
  });

  it("transforma a data em um Date", () => {
    const dados = esquemaNovaTarefa.parse({
      title: "Entregar",
      dueDate: "2026-12-01",
    });
    expect(dados.dueDate).toEqual(new Date("2026-12-01"));
  });

  it("descarta campos que não são da tarefa, como userId", () => {
    const dados = esquemaNovaTarefa.parse({
      title: "Invasão",
      userId: "id-de-outra-pessoa",
    });
    expect(dados).not.toHaveProperty("userId");
  });
});
\end{Verbatim}
\end{codigo}

O \texttt{describe} agrupa testes relacionados, e o \texttt{it} define um teste: um nome, que diz o comportamento esperado, e uma função que o confere. Dentro dela, o \texttt{expect} recebe o valor obtido, e o \textbf{\textit{matcher}} diz o que se espera dele: \texttt{toBe} compara valores simples, como números e textos; \texttt{toEqual} compara objetos campo a campo; \texttt{not.toHaveProperty} confere que um campo não existe. Se alguma expectativa não se cumpre, o \texttt{expect} lança um erro, e o teste falha.

Os testes seguem uma estrutura conhecida como \textbf{AAA}: \textit{Arrange}, preparar os dados; \textit{Act}, executar o que está sendo testado; \textit{Assert}, conferir o resultado. Nos testes curtos, as três partes cabem em duas linhas; nos maiores, uma linha em branco entre elas ajuda a leitura. Repare também no último teste: ele confere a proteção contra a atribuição em massa do Capítulo 6, uma regra de segurança que nenhuma revisão futura do esquema pode quebrar sem que alguém perceba.

\begin{codigo}
\begin{Verbatim}[breaklines=true, breaksymbolleft={}, fontsize=\footnotesize]
describe("esquemaFiltros", () => {
  it("preenche os valores padrão", () => {
    expect(esquemaFiltros.parse({})).toEqual({
      sort: "-createdAt",
      page: 1,
      limit: 20,
    });
  });

  it("converte os textos da query string em números", () => {
    const filtros = esquemaFiltros.parse({
      page: "2",
      limit: "5",
    });
    expect(filtros.page).toBe(2);
    expect(filtros.limit).toBe(5);
  });

  it.each([
    { query: { page: "0" }, campo: "page" },
    { query: { page: "abc" }, campo: "page" },
    { query: { limit: "500" }, campo: "limit" },
    { query: { sort: "nome" }, campo: "sort" },
    { query: { status: "FEITO" }, campo: "status" },
  ])("recusa $query", ({ query, campo }) => {
    const resultado = esquemaFiltros.safeParse(query);
    expect(resultado.success).toBe(false);
    expect(resultado.error?.issues[0].path).toEqual([campo]);
  });
});

describe("esquemaCadastro", () => {
  it("guarda o e-mail sem espaços e em minúsculas", () => {
    const dados = esquemaCadastro.parse({
      name: "Ana",
      email: " Ana@Mail.COM ",
      password: "senha12345",
    });
    expect(dados.email).toBe("ana@mail.com");
  });

  it("recusa senha de mais de 72 bytes", () => {
    const resultado = esquemaCadastro.safeParse({
      name: "Ana",
      email: "ana@mail.com",
      password: "á".repeat(40), // 40 letras, 80 bytes
    });
    expect(resultado.success).toBe(false);
  });
});
\end{Verbatim}
\end{codigo}

O \texttt{it.each} roda o mesmo teste com vários dados, um para cada item da lista. O \texttt{\$query} no nome é trocado pelo valor de cada caso, e cada um aparece como um teste separado no resultado: ``recusa \{ page: '0' \}'', ``recusa \{ limit: '500' \}''. Escrever cinco testes quase iguais seria cansativo e, pior, convidaria a testar menos casos. O teste de senha usa a letra ``á'', que ocupa 2 bytes: 40 letras passam do limite de 72 bytes do Capítulo 9, embora sejam menos de 72 caracteres.

\subsection{11.6 Testando os erros e usando espiões}

O módulo de erros do Capítulo 8 é outro bom candidato a testes de unidade:

\begin{codigo}
\begin{Verbatim}[breaklines=true, breaksymbolleft={}, fontsize=\footnotesize]
// testes/unidade/erros.test.ts -- o formato dos erros
import { describe, expect, it, vi } from "vitest";
import { ErroApi, problema, respostaDeErro } from "@/lib/erros";
import { validar } from "@/lib/validar";
import { esquemaNovaTarefa } from "@/lib/esquemas/tarefas";

describe("problema", () => {
  it("monta uma resposta no formato da RFC 9457", async () => {
    const resposta = problema(404, "Tarefa não encontrada.");

    expect(resposta.status).toBe(404);
    expect(resposta.headers.get("Content-Type")).toBe(
      "application/problem+json",
    );
    expect(await resposta.json()).toEqual({
      title: "Não encontrado",
      status: 404,
      detail: "Tarefa não encontrada.",
    });
  });

  it("acrescenta o WWW-Authenticate só no 401", () => {
    expect(
      problema(401, "x").headers.get("WWW-Authenticate"),
    ).toBe("Bearer");
    expect(problema(403, "x").headers.has("WWW-Authenticate"))
      .toBe(false);
  });
});

describe("validar", () => {
  it("devolve os dados quando são válidos", () => {
    expect(validar(esquemaNovaTarefa, { title: "Ok" }))
      .toEqual({ title: "Ok" });
  });

  it("lança um ErroApi 400 com os erros por campo", () => {
    expect(() => validar(esquemaNovaTarefa, { title: "" }))
      .toThrow(expect.objectContaining({
        status: 400,
        campos: { title: ["O título é obrigatório."] },
      }));
  });
});
\end{Verbatim}
\end{codigo}

Para conferir que uma função \textbf{lança} um erro, ela não pode ser chamada diretamente, porque o erro interromperia o teste. Ela vai dentro de uma função, \texttt{() => validar(...)}, que o \texttt{expect} chama e observa. O \texttt{toThrow} aceita o tipo de erro esperado ou, como aqui, um \texttt{expect.objectContaining}, que confere só os campos indicados do erro lançado: o \texttt{status} e os erros por campo.

O teste mais interessante é o do erro inesperado. A \texttt{respostaDeErro} escreve no \textit{log} com o \texttt{console.error}, e o teste precisa conferir \textbf{o que} foi escrito. Para isso servem os \textbf{espiões}:

\begin{codigo}
\begin{Verbatim}[breaklines=true, breaksymbolleft={}, fontsize=\footnotesize]
describe("respostaDeErro", () => {
  it("transforma um ErroApi na resposta que ele descreve", () => {
    const resposta = respostaDeErro(
      new ErroApi(409, "Conflito."),
    );
    expect(resposta.status).toBe(409);
  });

  it("esconde o erro inesperado e o registra", async () => {
    // troca o console.error por uma função espiã, que só anota
    const log = vi.spyOn(console, "error")
      .mockImplementation(() => {});

    const resposta = respostaDeErro(
      new Error("falha no banco: senha=123"),
    );
    const corpo = await resposta.json();

    expect(resposta.status).toBe(500);
    expect(JSON.stringify(corpo)).not.toContain("senha=123");

    // o mesmo identificador vai para a resposta e para o log
    const id = corpo.instance.replace("urn:uuid:", "");
    expect(log).toHaveBeenCalledWith(
      `[erro ${id}]`,
      expect.any(Error),
    );
    log.mockRestore();
  });
});
\end{Verbatim}
\end{codigo}

O \texttt{vi.spyOn(console, "error")} troca o \texttt{console.error} por um espião, que anota cada chamada, com os argumentos. O \texttt{mockImplementation(() => \{\})} faz o espião não escrever nada, para o erro de propósito não poluir o resultado dos testes. O teste confere as duas regras do Capítulo 8: a mensagem interna, com a ``senha'', não vai para o cliente, e o identificador da resposta é o mesmo do \textit{log}. O \texttt{expect.any(Error)} aceita qualquer objeto do tipo \texttt{Error}. No fim, o \texttt{mockRestore} devolve o \texttt{console.error} original, para não afetar os outros testes.

\subsection{11.7 Controlando o relógio}

Como testar que um \textit{token} vence depois de uma hora, sem esperar uma hora? O Vitest permite trocar o relógio do JavaScript por um relógio falso:

\begin{codigo}
\begin{Verbatim}[breaklines=true, breaksymbolleft={}, fontsize=\footnotesize]
// testes/unidade/token.test.ts -- os tokens JWT
import { afterEach, describe, expect, it, vi } from "vitest";
import { conferirToken, criarToken } from "@/lib/token";

describe("tokens", () => {
  afterEach(() => {
    vi.useRealTimers();
  });

  it("aceita o token que criou e devolve o id", async () => {
    const { token, expiresIn } = await criarToken("usuario-1");
    expect(expiresIn).toBe(3600);
    expect(await conferirToken(token)).toBe("usuario-1");
  });

  it("recusa um token com a carga alterada", async () => {
    const { token } = await criarToken("usuario-1");
    const [cabecalho, , assinatura] = token.split(".");
    const outraCarga = Buffer.from(
      JSON.stringify({ sub: "usuario-2" }),
    ).toString("base64url");

    const falso = [cabecalho, outraCarga, assinatura].join(".");
    expect(await conferirToken(falso)).toBeNull();
  });

  it("recusa textos que não são um JWT", async () => {
    expect(await conferirToken("abc")).toBeNull();
    expect(await conferirToken("")).toBeNull();
  });

  it("recusa o token depois de uma hora", async () => {
    // relógio falso: o teste não precisa esperar uma hora
    vi.useFakeTimers({ toFake: ["Date"] });
    vi.setSystemTime(new Date("2026-10-17T10:00:00Z"));
    const { token } = await criarToken("usuario-1");

    vi.setSystemTime(new Date("2026-10-17T10:59:00Z"));
    expect(await conferirToken(token)).toBe("usuario-1");

    vi.setSystemTime(new Date("2026-10-17T11:01:00Z"));
    expect(await conferirToken(token)).toBeNull();
  });
});
\end{Verbatim}
\end{codigo}

O \texttt{vi.useFakeTimers(\{ toFake: ["Date"] \})} troca só o \texttt{Date}, que é o que a \texttt{jose} consulta para saber a hora atual; as outras funções de tempo continuam verdadeiras. O \texttt{vi.setSystemTime} acerta o relógio falso em qualquer data: o \textit{token} é criado às 10h, ainda vale às 10h59 e é recusado às 11h01. O \texttt{afterEach} roda depois de cada teste e devolve o relógio verdadeiro, mesmo que o teste falhe no meio.

O segundo teste refaz o ataque do Capítulo 9: troca a carga do \textit{token}, mantendo a assinatura original. Os testes de segurança são dos mais valiosos, porque o defeito que eles pegam não aparece no uso normal: ninguém percebe que a API aceita um \textit{token} falso até que alguém o envie.

\subsection{11.8 Testes de integração: os serviços com o banco}

Os serviços do Capítulo 7 são o coração da API, e o que eles fazem depende do banco: os filtros, a ordenação e, principalmente, o \texttt{userId} em cada consulta. Testá-los com um banco falso esconderia justamente o que pode dar errado. Os testes de integração usam o banco dos testes:

\begin{codigo}
\begin{Verbatim}[breaklines=true, breaksymbolleft={}, fontsize=\footnotesize]
// testes/integracao/usuarios.test.ts -- cadastro e login
import { beforeEach, describe, expect, it } from "vitest";
import { db } from "@/lib/db";
import {
  autenticar,
  cadastrarUsuario,
} from "@/lib/servicos/usuarios";
import { limparBanco } from "../apoio";

const ana = {
  name: "Ana",
  email: "ana@teste.com",
  password: "senha12345",
};

describe("serviço de usuários", () => {
  beforeEach(limparBanco);

  it("cadastra sem devolver a senha; guarda o hash", async () => {
    const usuario = await cadastrarUsuario(ana);
    expect(usuario).not.toHaveProperty("password");

    const noBanco = await db.user.findUniqueOrThrow({
      where: { email: ana.email },
    });
    expect(noBanco.password).not.toBe(ana.password);
    expect(noBanco.password).toMatch(/^\$2b\$10\$/);
  });

  it("devolve null quando o e-mail já existe", async () => {
    await cadastrarUsuario(ana);
    expect(await cadastrarUsuario(ana)).toBeNull();
  });

  it("autentica só com o e-mail e a senha certos", async () => {
    const usuario = await cadastrarUsuario(ana);

    expect(await autenticar(ana)).toBe(usuario?.id);
    expect(
      await autenticar({ ...ana, password: "errada123" }),
    ).toBeNull();
    expect(
      await autenticar({ ...ana, email: "nao@existe.com" }),
    ).toBeNull();
  });
});
\end{Verbatim}
\end{codigo}

O \texttt{beforeEach(limparBanco)} roda a limpeza antes de cada teste do grupo. Assim, cada teste começa do zero e pode ser rodado sozinho, em qualquer ordem, com o mesmo resultado. O \texttt{toMatch} confere um texto contra uma expressão regular: o que está guardado no banco tem de começar com \texttt{\$2b\$10\$}, o início de um \textit{hash} do bcrypt com custo 10. Agora, as tarefas:

\begin{codigo}
\begin{Verbatim}[breaklines=true, breaksymbolleft={}, fontsize=\footnotesize]
// testes/integracao/tarefas.test.ts -- o serviço de tarefas
import { beforeEach, describe, expect, it } from "vitest";
import { esquemaFiltros } from "@/lib/esquemas/tarefas";
import {
  buscarTarefa,
  criarTarefa,
  editarTarefa,
  excluirTarefa,
  listarTarefas,
} from "@/lib/servicos/tarefas";
import { criarPessoa, limparBanco } from "../apoio";

// Os filtros completos, com os padrões, a partir da query string
const filtros = (query: Record<string, string> = {}) =>
  esquemaFiltros.parse(query);

describe("serviço de tarefas", () => {
  beforeEach(limparBanco);

  it("cria uma tarefa com os padrões do banco", async () => {
    const ana = await criarPessoa();
    const tarefa = await criarTarefa(ana.id, {
      title: "Estudar",
    });

    expect(tarefa).toMatchObject({
      title: "Estudar",
      status: "PENDING",
      priority: "MEDIUM",
      dueDate: null,
    });
    expect(tarefa).not.toHaveProperty("userId");
  });

  it("não deixa mexer na tarefa de outra pessoa", async () => {
    const ana = await criarPessoa("Ana");
    const bia = await criarPessoa("Bia");
    const tarefa = await criarTarefa(ana.id, { title: "Da Ana" });

    expect(await buscarTarefa(bia.id, tarefa.id)).toBeNull();
    expect(
      await editarTarefa(bia.id, tarefa.id, { title: "Outro" }),
    ).toBeNull();
    expect(await excluirTarefa(bia.id, tarefa.id)).toBe(false);

    // e a tarefa continua intacta para a dona
    const daAna = await buscarTarefa(ana.id, tarefa.id);
    expect(daAna?.title).toBe("Da Ana");
  });
\end{Verbatim}
\end{codigo}

O \texttt{toMatchObject} confere só os campos indicados e ignora os outros, como o \texttt{id} e as datas, que mudam a cada execução. O segundo teste é o mais importante do arquivo. Ele confere a regra de segurança do Capítulo 2, que o OWASP chama de BOLA (Capítulo 12): a Bia não consegue ver, editar nem excluir a tarefa da Ana, e a tarefa continua intacta depois das tentativas.

\begin{codigo}
\begin{Verbatim}[breaklines=true, breaksymbolleft={}, fontsize=\footnotesize]
  it("filtra, busca e pagina a lista", async () => {
    const ana = await criarPessoa();
    for (let i = 1; i <= 5; i++) {
      await criarTarefa(ana.id, { title: `Tarefa ${i}` });
    }
    await criarTarefa(ana.id, {
      title: "Relatório final",
      status: "DONE",
    });

    const feitas = await listarTarefas(
      ana.id,
      filtros({ status: "DONE" }),
    );
    expect(feitas.data.map((t) => t.title))
      .toEqual(["Relatório final"]);

    const busca = await listarTarefas(
      ana.id,
      filtros({ search: "RELAT" }),
    );
    expect(busca.meta.total).toBe(1);

    const pagina2 = await listarTarefas(
      ana.id,
      filtros({ limit: "4", page: "2" }),
    );
    expect(pagina2.data).toHaveLength(2);
    expect(pagina2.meta).toEqual({
      page: 2,
      limit: 4,
      total: 6,
      totalPages: 2,
    });
  });

  it("ordena pela data, com as sem data no fim", async () => {
    const ana = await criarPessoa();
    await criarTarefa(ana.id, { title: "Sem data" });
    await criarTarefa(ana.id, {
      title: "Dezembro",
      dueDate: new Date("2026-12-01"),
    });
    await criarTarefa(ana.id, {
      title: "Novembro",
      dueDate: new Date("2026-11-05"),
    });

    for (const sort of ["dueDate", "-dueDate"]) {
      const lista = await listarTarefas(
        ana.id,
        filtros({ sort }),
      );
      expect(lista.data.at(-1)?.title).toBe("Sem data");
    }
  });
});
\end{Verbatim}
\end{codigo}

A função \texttt{filtros} passa a \textit{query string} pelo \texttt{esquemaFiltros}, para obter o objeto completo, com os padrões, exatamente como a rota faz. O teste da ordenação confere, nas duas direções, que as tarefas sem data ficam no fim, o detalhe do \texttt{nulls: "last"} do Capítulo 7. O \texttt{at(-1)} pega o último item de uma lista.

\subsection{11.9 Testando as rotas sem o servidor}

Uma rota do Next.js é uma função exportada: dá para importá-la e chamá-la no teste, com uma requisição montada à mão, sem subir o servidor. O obstáculo é a \texttt{usuarioAtual}, que lê os cabeçalhos com o \texttt{headers()} do Next.js, e essa função só funciona dentro do servidor. A solução é um \textbf{\textit{mock}}: uma versão falsa de um módulo, usada só nos testes.

\begin{codigo}
\begin{Verbatim}[breaklines=true, breaksymbolleft={}, fontsize=\footnotesize]
// testes/integracao/rotas.test.ts -- as rotas, sem o servidor
import { beforeEach, describe, expect, it, vi } from "vitest";
import { NextRequest } from "next/server";
import { Ajv2020 } from "ajv/dist/2020.js";
import * as colecao from "@/app/api/tasks/route";
import * as umaTarefa from "@/app/api/tasks/[id]/route";
import { documento } from "@/lib/openapi";
import { criarToken } from "@/lib/token";
import { criarPessoa, limparBanco } from "../apoio";

// A usuarioAtual lê os cabeçalhos com o headers() do Next.js,
// que só funciona com o servidor no ar. No teste, o módulo é
// trocado por um falso (mock), que devolve estes cabeçalhos.
const cabecalhos = vi.hoisted(() => new Headers());
vi.mock("next/headers", () => ({
  headers: async () => cabecalhos,
}));

async function entrarComo(userId: string) {
  const { token } = await criarToken(userId);
  cabecalhos.set("Authorization", `Bearer ${token}`);
}

const requisicao = (url: string, corpo?: unknown) =>
  new NextRequest(`http://localhost${url}`, {
    method: corpo === undefined ? "GET" : "POST",
    body: corpo === undefined ? undefined : JSON.stringify(corpo),
  });

const contexto = (id: string) => ({
  params: Promise.resolve({ id }),
});
\end{Verbatim}
\end{codigo}

O \texttt{vi.mock("next/headers", ...)} faz todo \texttt{import} desse módulo, inclusive o que está dentro da \texttt{usuarioAtual}, receber o objeto devolvido pela função, em que o \texttt{headers} devolve os cabeçalhos do teste. O Vitest move o \texttt{vi.mock} para o começo do arquivo, antes dos \texttt{import}, porque o módulo falso precisa existir antes que alguém o importe. Por isso a variável \texttt{cabecalhos} é criada com o \texttt{vi.hoisted}, que também é movido para o começo; uma variável comum ainda não existiria quando o \textit{mock} fosse criado. A \texttt{entrarComo} cria um \textit{token} de verdade e o coloca no cabeçalho, como faria um cliente depois do \textit{login}.

\begin{codigo}
\begin{Verbatim}[breaklines=true, breaksymbolleft={}, fontsize=\footnotesize]
describe("rotas de tarefas", () => {
  beforeEach(async () => {
    await limparBanco();
    cabecalhos.delete("Authorization");
  });

  it("responde 401 sem token, antes de validar", async () => {
    const resposta = await colecao.POST(
      requisicao("/api/tasks", { title: "" }),
      {},
    );
    expect(resposta.status).toBe(401);
    expect(resposta.headers.get("WWW-Authenticate"))
      .toBe("Bearer");
  });

  it("cria com 201 e Location, e a tarefa existe", async () => {
    const ana = await criarPessoa();
    await entrarComo(ana.id);

    const criada = await colecao.POST(
      requisicao("/api/tasks", { title: "Pela rota" }),
      {},
    );
    expect(criada.status).toBe(201);
    const { id } = await criada.json();
    expect(criada.headers.get("Location"))
      .toBe(`/api/tasks/${id}`);

    const lida = await umaTarefa.GET(
      requisicao(`/api/tasks/${id}`),
      contexto(id),
    );
    expect(lida.status).toBe(200);
  });

  it("responde 400 no formato Problem Details", async () => {
    const ana = await criarPessoa();
    await entrarComo(ana.id);

    const resposta = await colecao.POST(
      requisicao("/api/tasks", {
        title: "",
        priority: "URGENTE",
      }),
      {},
    );
    expect(resposta.status).toBe(400);
    const corpo = await resposta.json();
    expect(Object.keys(corpo.errors))
      .toEqual(["title", "priority"]);
  });
});
\end{Verbatim}
\end{codigo}

Os testes chamam o \texttt{POST} e o \texttt{GET} exportados pelas rotas, que já vêm envolvidos pela \texttt{comErros}; por isso a resposta de erro é a mesma que o cliente receberia. O \texttt{\{\}} é o contexto da rota da coleção, que não tem parâmetros. O primeiro teste confere a ordem combinada no Capítulo 8: sem \textit{token}, a resposta é \texttt{401}, mesmo com um título vazio, que daria \texttt{400}.

Repare no que ficou de fora do \textit{mock}: o \textit{token}, a validação, o serviço e o banco são todos verdadeiros. Só foi trocado o que não funciona fora do servidor. Quanto mais coisas um teste troca por versões falsas, menos ele prova sobre o sistema de verdade.

\subsection{11.10 O teste de contrato}

O Capítulo 10 deixou um risco anotado: os esquemas das \textbf{respostas} no documento OpenAPI foram escritos à mão, e podem deixar de corresponder ao que a API devolve. Um \textbf{teste de contrato} confere a resposta real contra o documento:

\begin{codigo}
\begin{Verbatim}[breaklines=true, breaksymbolleft={}, fontsize=\footnotesize]
describe("contrato: as respostas seguem o OpenAPI", () => {
  // O Ajv confere dados contra um JSON Schema; os esquemas vêm
  // dos components do documento do Capítulo 10
  const ajv = new Ajv2020({ validateFormats: false });
  for (const [nome, esquema] of Object.entries(
    documento.components.schemas,
  )) {
    ajv.addSchema(esquema, `#/components/schemas/${nome}`);
  }

  beforeEach(async () => {
    await limparBanco();
    cabecalhos.delete("Authorization");
  });

  it("GET /api/tasks devolve uma ListaDeTarefas", async () => {
    const ana = await criarPessoa();
    await entrarComo(ana.id);
    await colecao.POST(
      requisicao("/api/tasks", {
        title: "Com data",
        dueDate: "2026-12-01",
      }),
      {},
    );

    const resposta = await colecao.GET(
      requisicao("/api/tasks"),
      {},
    );
    const confere = ajv.getSchema(
      "#/components/schemas/ListaDeTarefas",
    );
    const valido = confere?.(await resposta.json());
    expect(confere?.errors ?? []).toEqual([]);
    expect(valido).toBe(true);
  });
});
\end{Verbatim}
\end{codigo}

O \textbf{Ajv} é a biblioteca mais usada para conferir dados contra um JSON Schema. A classe \texttt{Ajv2020} entende a versão 2020-12 do JSON Schema, a do OpenAPI 3.1. O laço registra cada esquema do documento com o mesmo endereço usado nos \texttt{\$ref}, como \texttt{\#/components/schemas/Tarefa}; assim, quando a \texttt{ListaDeTarefas} aponta para a \texttt{Tarefa}, o Ajv sabe onde encontrá-la. O \texttt{getSchema} devolve uma função que confere os dados e deixa os problemas encontrados em \texttt{errors}. Comparar o \texttt{errors} com uma lista vazia, antes de conferir o \texttt{true}, faz o Vitest mostrar quais campos não bateram, se o teste falhar. A opção \texttt{validateFormats: false} dispensa o Ajv de conferir formatos como \texttt{date-time}: as datas já são conferidas pelos padrões (\texttt{pattern}) que o Zod incluiu nos esquemas.

Se alguém acrescentar um campo ao \texttt{camposTarefa} e esquecer o documento, este teste ainda passa, porque o JSON Schema aceita campos extras por padrão. Mas se um campo mudar de tipo, ou sumir da resposta, ele falha, e a documentação errada é descoberta antes de chegar a quem usa a API.

\subsection{11.11 Rodando os testes}

Com o banco dos testes criado, um comando roda tudo:

\begin{codigo}
\begin{Verbatim}[breaklines=true, breaksymbolleft={}, fontsize=\footnotesize]
$ npx vitest run

 RUN  v5.0.5

1 migration found in prisma/migrations
No pending migrations to apply.

 Test Files  6 passed (6)
      Tests  34 passed (34)
   Duration  10.85s
\end{Verbatim}
\end{codigo}

Seis arquivos e 34 testes, todos aprovados. Para ver como aparece uma falha, a busca de uma tarefa, no serviço do Capítulo 7, foi alterada de propósito: o \texttt{where: \{ id, userId \}} virou \texttt{where: \{ id \}}, um esquecimento que deixaria qualquer pessoa ler a tarefa de qualquer outra. O resultado, com os símbolos do Vitest trocados por \texttt{>} e \texttt{X}, alguns espaços reduzidos e a tarefa resumida:

\begin{codigo}
\begin{Verbatim}[breaklines=true, breaksymbolleft={}, fontsize=\footnotesize]
 > testes/integracao/tarefas.test.ts (4 tests | 1 failed)
   > serviço de tarefas (4)
     X não deixa mexer na tarefa de outra pessoa

 FAIL  testes/integracao/tarefas.test.ts > serviço de tarefas
       > não deixa mexer na tarefa de outra pessoa
AssertionError: expected { ...(8) } to be null
- Expected:
null
+ Received:
{
  "id": "cmvcv6xxo00040ov9b7h5jp2k",
  "title": "Da Ana",
  ...
}
 > testes/integracao/tarefas.test.ts:40:51
38| const tarefa = await criarTarefa(ana.id, { title: "Da Ana" });
39|
40| expect(await buscarTarefa(bia.id, tarefa.id)).toBeNull();
\end{Verbatim}
\end{codigo}

O Vitest diz qual teste falhou, o que era esperado (\texttt{null}), o que veio (a tarefa da Ana) e em qual linha. O nome do teste já explica a gravidade do problema. Com o \texttt{userId} de volta, os 34 testes voltaram a passar.

A \textbf{cobertura} mede quanto do código os testes executaram. O \texttt{npm run cobertura} mostra uma tabela; este é um trecho do resultado, com as linhas de cada pasta:

\begin{codigo}
\begin{Verbatim}[breaklines=true, breaksymbolleft={}, fontsize=\footnotesize]
File               | % Stmts | % Branch | % Funcs | % Lines
-------------------|---------|----------|---------|--------
All files          |   67.51 |    67.69 |   58.97 |   69.33
 .../auth/register |       0 |        0 |       0 |       0
 app/api/tasks     |   91.66 |      100 |   66.66 |     100
 ...api/tasks/[id] |      40 |    16.66 |      20 |    42.1
  erros.ts         |     100 |    85.71 |     100 |     100
  token.ts         |    90.9 |    66.66 |     100 |    90.9
  xml.ts           |       0 |        0 |       0 |       0
 lib/servicos      |   93.54 |       90 |     100 |   93.54
\end{Verbatim}
\end{codigo}

As colunas contam instruções (\textit{Stmts}), caminhos de \texttt{if} e \texttt{?:} (\textit{Branch}), funções e linhas. A tabela mostra o que ainda não foi testado: as rotas de cadastro e de \textit{login}, e o \texttt{PATCH} e o \texttt{DELETE} da rota de uma tarefa, cujos serviços foram testados, mas as rotas, não. Ela também mostra o \texttt{lib/xml.ts} com zero: desde o Capítulo 7, nenhuma rota usa a resposta em XML. A cobertura encontra também código que ninguém usa.

Use a cobertura para descobrir o que falta testar, e não como meta. Cem por cento de cobertura não garante nada: um teste pode executar uma linha sem conferir o resultado dela. E correr atrás dos últimos pontos leva a testes que conferem detalhes sem importância. O que importa é que as regras de negócio e de segurança estejam cobertas, como estão, no TaskFlow, os serviços, os erros e os \textit{tokens}.

\subsection{11.12 Bons testes}

Alguns hábitos separam os testes que ajudam dos que atrapalham:

\begin{itemize}
  \item \textbf{Independentes.} Cada teste prepara o que precisa e não depende de outro ter rodado antes. Por isso o banco é limpo antes de cada um.
  \item \textbf{Determinísticos.} O mesmo código dá sempre o mesmo resultado. Datas, horas e valores aleatórios são as causas mais comuns de testes que falham ``de vez em quando''; controle-os, como o relógio falso fez com o \textit{token}.
  \item \textbf{Testam o comportamento, e não a implementação.} O teste do isolamento confere que a Bia não vê a tarefa da Ana, e não que a consulta tem um \texttt{userId}. Se a consulta for reescrita de outro jeito, o teste continua valendo.
  \item \textbf{Nomes que explicam.} O nome diz a regra: ``recusa o \textit{token} depois de uma hora''. Quando ele falhar, a mensagem já diz o que quebrou.
  \item \textbf{Testam os erros.} O caminho feliz é o que menos quebra. Dados inválidos, \textit{tokens} falsos e recursos de outra pessoa são onde estão os defeitos e as falhas de segurança.
  \item \textbf{Viu falhar, confiou.} Um teste que nunca falhou pode não estar testando nada. Ao escrever um teste novo, quebre o código de propósito, como foi feito com o \texttt{userId}, e confira que ele falha.
\end{itemize}

Os três tipos de teste do TaskFlow se completam: o Vitest confere as regras, por dentro, em segundos; a coleção do Postman, rodada pelo Newman, confere a API de fora, com o servidor no ar, como o cliente a vê. No Capítulo 12, os dois passam a rodar sozinhos a cada envio de código ao GitHub, e uma falha impede a publicação. O projeto fica assim:

\begin{codigo}
\begin{Verbatim}[breaklines=true, breaksymbolleft={}, fontsize=\footnotesize]
.env.test                  o banco dos testes (fora do Git)
vitest.config.mts          a configuração do Vitest
testes/
  preparar-banco.ts        a trava e as migrações
  apoio.ts                 limparBanco, criarPessoa
  unidade/
    esquemas.test.ts       o que a API aceita
    erros.test.ts          o formato dos erros, espiões
    token.test.ts          os tokens, relógio falso
  integracao/
    usuarios.test.ts       cadastro e login, com o banco
    tarefas.test.ts        serviço de tarefas, isolamento
    rotas.test.ts          rotas com mock e contrato
\end{Verbatim}
\end{codigo}

\subsection{Erros Comuns}

\begin{itemize}
  \item \textbf{Rodar os testes no banco de desenvolvimento ou de produção.} Os testes apagam dados. Use um banco só para eles e uma trava que confira isso antes de começar.
  \item \textbf{Testes que dependem da ordem.} Um teste que só passa se outro rodou antes quebra ao ser rodado sozinho. Prepare tudo dentro de cada teste ou no \texttt{beforeEach}.
  \item \textbf{Rodar em paralelo testes que dividem o banco.} Um apaga os dados do outro, e as falhas parecem aleatórias. Use o \texttt{fileParallelism: false}, ou um banco por processo.
  \item \textbf{Chamar diretamente a função que deve lançar o erro.} O erro interrompe o teste. Passe uma função ao \texttt{expect}: \texttt{expect(() => validar(...)).toThrow()}.
  \item \textbf{Esquecer o \texttt{await}.} Sem ele, o \texttt{expect} compara uma \textit{Promise}, e não o resultado. Em funções \texttt{async}, use \texttt{await} em toda chamada que devolve uma \textit{Promise}.
  \item \textbf{Trocar tudo por \textit{mocks}.} Um teste em que o banco, o serviço e a validação são falsos só prova que os \textit{mocks} funcionam. Troque apenas o que não funciona no ambiente de testes.
  \item \textbf{Não restaurar espiões e relógios.} Um \texttt{console.error} silenciado ou um relógio falso esquecido afetam os testes seguintes. Use \texttt{mockRestore} e \texttt{useRealTimers}.
  \item \textbf{Perseguir 100\% de cobertura.} A cobertura mostra o que não foi executado, não o que foi bem testado. Priorize as regras de negócio e de segurança.
\end{itemize}
